\section{Related Work}
\label{sec:related_work}

\noindent \textbf{Generalizable 3D Gaussian Splatting.}
Single-scene 3D Gaussian Splatting methods~\cite{kerbl20233d, fan2024instantsplat, yan2024street, niemeyer2025radsplat, feng2025flashgs} enable more efficient rendering than neural fields and volume rendering methods~\cite{lombardi2019neural, sitzmann2020implicit, mildenhall2021nerf, xie2022neural} thanks to their rasterization-friendly formulation, but they still suffer from long optimization time and limited generalization.
%
To address this issue, generalizable 3D Gaussian Splatting methods~\cite{charatan2024pixelsplat, tang2024hisplat, liu2025monosplat, xu2025depthsplat} predict all Gaussian parameters in a single feed-forward pass, enabling fast reconstruction and novel view synthesis for unseen scenes.
%
Early works~\cite{charatan2024pixelsplat, szymanowicz2024splatter} directly regressed Gaussian parameters from image features using feed-forward networks.
%
With the introduction of cost-volume construction into generalizable 3DGS, subsequent methods~\cite{chen2024mvsplat, liu2024mvsgaussian} exploited cross-view geometric cues to improve depth estimation and simplify Gaussian parameter prediction.
%
More recently, methods~\cite{xu2025depthsplat, liu2025monosplat} further improved performance by incorporating pre-trained monocular depth models.
%
However, existing methods still rely on a single warp for depth estimation, which limits their ability to fully exploit cross-view feature cues.
%
In contrast, IDESplat integrates feature similarity from cascaded warps to produce more reliable and refined depth maps for Gaussian mean prediction.


\noindent \textbf{Iteration-based Optimization Methods.}
%
%
%
\noindent Iteration-based optimization methods~\cite{adler2017solving, li2018deepim} are widely used in various tasks~\cite{he2024lotus, zhu2024addressing, chen2024virtual} due to their ability to progressively enhance and refine the feature learning process.
%
For instance, in optical flow estimation~\cite{teed2020raft, hui2018liteflownet,  yang2019volumetric, sun2018pwc, ilg2017flownet}, iterative optimization is commonly used to build coarse-to-fine pyramidal features by stacking multiple feature units, resulting in more accurate and stable flow estimation.
%
In the field of depth estimation, some works~\cite{lipson2021raft, wang2022itermvs, wang2022efficient, ma2022multiview, xu2023iterative, xu2025igev++, chen2024mocha} have iteratively retrieved multi-view correlation features using structures like GRUs and LSTMs with shared parameters to update the disparity field, yielding better depth estimation results.
%
Similarly, in monocular depth estimation~\cite{zuo2024ogni, shao2023iebins, fu2018deep}, iterative methods are often used to gradually refine the depth search range, leading to more stable and accurate predictions.
%
Additionally, in the 3D Gaussian Splatting (3DGS) domain, recent methods~\cite{wenlife, zhang2025transplat, xu2025resplat, chen2024g3r, harrison2022closer} have also attempted to design iterative optimization processes to reduce the difficulty of 3D Gaussian reconstruction tasks.
%
%
Unlike existing iterative optimization methods, we propose a novel Depth Probability Estimation Unit that integrates the similarity results of multiple warps in a multiplicative manner.
%
This approach produces more reliable and refined depth maps, while progressively enhancing feature resolution and refining the depth candidate range throughout the iterative process.




